(1) 小角度时水平位移为 \(x_i\simeq\ell\theta_i\)，动能 \(T=m\ell^2(\dot\theta_1^2+\dot\theta_2^2)/2\)，势能 \(V=mg\ell(\theta_1^2+\theta_2^2)/2+k\ell^2(\theta_1-\theta_2)^2/2\)。Lagrange方程是 \(\ddot\theta_i+(g/\ell)\theta_i+(k/m)(\theta_i-\theta_j)=0\)。

(2) 令 \(q_+=(\theta_1+\theta_2)/2\)、\(q_-=(\theta_1-\theta_2)/2\)，两式分别变成 \(\ddot q_++\omega_+^2q_+=0\)、\(\ddot q_-+\omega_-^2q_-=0\)，其中 \(\boxed{\omega_+^2=g/\ell,\quad\omega_-^2=g/\ell+2k/m}\)。同相模式弹簧长度不变，因此仅有重力回复。

(3) 初态给 \(q_+(0)=q_-(0)=\Theta/2\) 且导数都为零。因此 \(\boxed{\theta_1(t)=\frac\Theta2(\cos\omega_+t+\cos\omega_-t)}\)，\(\boxed{\theta_2(t)=\frac\Theta2(\cos\omega_+t-\cos\omega_-t)}\)。使用和差化积可见快速振荡乘缓慢拍幅。

(4) 弱耦合时 \(\Delta\omega=\omega_--\omega_+\simeq k\ell/(m\sqrt{g\ell})\)。两摆的慢包络分别与 \(\cos(\Delta\omega t/2)\) 和 \(\sin(\Delta\omega t/2)\) 成比例，首个近乎完整的能量转移在 \(\boxed{t\simeq\pi/\Delta\omega}\)；严格能量转移会有快速相位修正，这里采用 \(\Delta\omega\ll\omega_+\) 的平均近似。
